\chapter{结论与展望}

\section{主要结论}

本文研究了金属热处理温度场的受控递归建模问题。模型以完整离散温度场为状态，以炉温为动作，并显式接收边界与材料参数。C45/AISI 1045 类中碳钢数值环境包含温变导热系数、温变比热和对流辐射边界。主要结论如下。

\begin{enumerate}
  \item 完整轨迹评价揭示了一步误差无法表达的长期行为。在相同数据、训练轮数和随机种子下，五步展开相对一步训练将 ID 轨迹 RMSE 降低 $4.2\%$，但控制 OOD 轨迹 RMSE 增加 $10.8\%$。因此，本文依据验证集轨迹误差选择五步设置，并把完整闭环轨迹作为模型选择与报告单位。
  \item 物理约束提高了分布外轨迹的能量一致性和尾部稳定性。验证选择的 $10^{-3}$ 权重在三随机种子上将控制 OOD RMSE 降低 $6.8\%$，物理残差降低 $44.8\%$。四参数组合 OOD 的滚动 RMSE 降低 $29.5\%$，最大绝对误差降低约 $49.3\%$。
  \item 动态对流辐射边界的瞬时温度响应只能识别 $\Heff=h+\varepsilon h_{\mathrm{rad}}$。该参数化将动态边界滚动 RMSE 从 $2.285\pm0.108$~\si{\degreeCelsius} 降至 $1.409\pm0.068$~\si{\degreeCelsius}。独立 BDF 求解器验证保持相同结论。
  \item 五点带噪温度观测配合集合滤波得到 \SI{1.008}{\degreeCelsius} 的全场 RMSE 和 \SI{0.231}{\degreeCelsius} 的中心 RMSE。ID 条件下闭环成功率为 $100\%$。严格参数 OOD 条件下，风险感知控制的综合目标平均配对差为 $-1.737$，95\% bootstrap 区间为 $[-3.291,-0.226]$。
  \item $\Heff$ 世界模型在独立 BDF 动力学上获得 55.4 倍完整轨迹推演加速，在候选动作规划中获得 63.4 倍加速。两项结果分别对应状态轨迹生成和滚动决策，说明该模型可以为高保真复核前的批量工艺筛选提供计算余量。
\end{enumerate}

上述结果支持一个明确判断。热处理世界模型的关键价值在于形成可递归、可观测、可用于动作比较的受控动力学表示。物理损失增强 OOD 稳定性，可辨识参数化减少动态边界歧义，状态估计器负责连接稀疏测量。三者共同决定模型能否进入闭环。

\section{研究范围}

本文证据来自一维对称平板数值实验。BDF 参考环境改变了离散网格和时间积分方法，但连续传热方程与材料物性函数保持一致。当前状态不含相变组织、热应力和变形。\SI{350}{\degreeCelsius} 目标跟踪用于检验控制链路，不构成完整的 C45 工业热处理制度。闭环实验尚未接入真实炉温、热电偶和执行器数据。因此，文中精度用于说明方法对当前热动力学环境的逼近和控制能力。

\section{后续工作}

后续研究可沿三个方向推进。第一，引入公开热电偶曲线或自建加热实验，对材料物性、表面发射率和传感器响应进行联合标定。第二，将状态模型扩展到二维截面和复杂几何，并比较局部卷积结构与神经算子。第三，建立材料 OOD 检测和混合规划器。世界模型负责快速筛选，高保真有限元模型复核少量高风险候选动作。该结构能够保留在线速度，也能提高复杂工况中的控制可信度。
